\section{Introduction}
\label{sec:intro}

Reinforcement learning from AI feedback (RLAIF) is increasingly used to train dialogue agents for
clinical and counselling tasks, where the reward is an LLM asked to fill in a validated
questionnaire about the conversation it just read. The appeal is obvious: the instrument is
published, the construct is one clinicians already argue about in the literature, and the grader
scales to tens of thousands of sessions. The risk is equally obvious, and has a name --- the
policy optimises the grader, not the construct.

In the setting we study, a 1B-parameter therapist trained against MI satisfaction and working
alliance rubrics acquires one specific shortcut. It learns to praise the client. Over ten
iterations, over-praise per session rises roughly eighteen-fold while every \emph{other}
MI-inconsistent behaviour --- confronting, warning, judging --- falls. The rubrics reward it
because they measure how the session felt to the client, and flattery is the cheapest way to move
that.

\paragraph{A principled intervention.} If the problem is that a turn is scored in isolation, score
the trajectory instead. Following the look-ahead formulation of preference-tree
optimisation, each candidate therapist turn is extended by $K$ further simulated turns --- the
client replies, the policy replies, and so on --- and the oracle grades the resulting continuation
rather than the turn. Praise that the conversation does not build on is supposed to stop looking
good. This is a reasonable, cheap, and general idea, and it is the kind of thing a practitioner
would reach for on being shown the flattery result.

\paragraph{What we find.} We run $K \in \{0, 5\}$ as a controlled pair and measure not the rubric
scores but the behaviours the rubrics summarise, persona-paired across 96 matched conversations at
each of eleven matched training iterations (0--10), under both the training oracle and a held-out
judge from a different model family.

Look-ahead does what it was supposed to do to the channel it targets
(\S\ref{sec:channel}). It does not reduce MI-inconsistent behaviour (\S\ref{sec:aggregate}). The
two arms commit the same number of MI violations per session; they commit different ones. What
$K{=}5$ saves in flattery it spends on unsolicited advice and directing the client --- and the
reward function the run was optimising is indifferent between the two.

\paragraph{Contributions.}
\begin{enumerate}
  \item A controlled test of trajectory-level reward as a reward-hacking intervention, with the
        candidate budget, optimiser, oracle and evaluation held fixed and $K$ as the only
        manipulation (\S\ref{sec:setup}).
  \item Evidence that the intervention closes its target channel at three measurement levels ---
        what the reward selects for, what the policy generates, and what the evaluation observes
        --- rather than only at the last of them (\S\ref{sec:mechanism}).
  \item Evidence that the aggregate the channel belongs to does not follow it: identical for eight
        iterations, and at the endpoint significant under the training oracle but null under the
        held-out judge, with every component effect large, opposed, and judge-robust
        (\S\ref{sec:aggregate}).
  \item Two methodological consequences: the aggregate has to be nominated \emph{before} the
        intervention, because any intervention looks successful when scored on the channel that
        motivated it; and an intervention that changes the \emph{composition} of a failure breaks
        the grader-invariance of aggregate comparisons, even when every component is
        grader-invariant (\S\ref{sec:discussion}).
\end{enumerate}

We do not claim look-ahead is useless, and we do not claim reward hacking is conserved as a law.
We claim that on this task, this reward, and this intervention, the quantity that moved robustly
was composition, that the quantity that moved unreliably was volume, and that the experiment which
would have reported a clean success is the one that measured only over-praise --- or only a rate.
